\section{Recommended Mathematical Background}

The main mathematical prerequisites are classical mechanics, multivariable
calculus, linear algebra, and ordinary differential equations. The reader should
be comfortable with gradients, divergence, curl, integration by parts, matrices,
eigenvectors, and basic variational ideas.

Partial differential equations appear throughout the course. Full mathematical
rigour about existence and uniqueness is not the main goal, but the reader
should understand that field equations require initial data, boundary data, and
conditions on the domain.

Some tensor notation and special relativity will also be used. The necessary
material is reviewed inside the notes, but the reader should expect the notation
to become more geometric as the course develops. Differential forms, Lie groups,
and spinors are introduced only as much as needed for classical field theory.

\begin{summarybox}
The mathematical attitude needed for these notes is more important than a long
list of prerequisites: follow definitions carefully, track assumptions, check
dimensions, and ask what each symbol represents physically.
\end{summarybox}
